% Balanced-growth calibration report -- TEMPLATE
% Numeric cells are \np until filled.  Figures live in \rundir.
% Compile: latexmk -pdf calibration_report.tex
\documentclass[11pt,a4paper]{article}
\usepackage[margin=1.9cm,top=1.6cm,bottom=1.6cm]{geometry}
\renewcommand{\topfraction}{.95}\renewcommand{\bottomfraction}{.95}
\renewcommand{\textfraction}{.05}\renewcommand{\floatpagefraction}{.9}
\usepackage{booktabs,graphicx,amsmath,amstext,siunitx,textcomp}
\newcommand{\EUR}[1]{\texteuro#1}
\usepackage[labelfont=bf,font=small,skip=4pt]{caption}

\sisetup{table-align-text-post=false}
\newcommand{\rundir}{../output/calibration_growth}
\newcommand{\fiscaldir}{../output/fiscal_2026-10-06_grid50}
% Fallbacks so the template compiles before any run exists; header.tex
% overrides them with the values of the run \rundir points at.
\newcommand{\runlabel}{\np}\newcommand{\gval}{\np}\newcommand{\nval}{\np}
\newcommand{\nTval}{\np}\newcommand{\rBval}{\np}\newcommand{\Gammaval}{\np}
\newcommand{\GammaTval}{\np}
\IfFileExists{\rundir/header.tex}{\input{\rundir/header.tex}}{}
\newcommand{\np}{\ensuremath{\text{--}}}   % empty cell; safe in S columns and math
\newcommand{\tabinput}[1]{\IfFileExists{\rundir/#1}{\input{\rundir/#1}}%
  {\figbox{run \texttt{fill\_report.py} to fill this table}{1.2cm}}}
\newcommand{\fiscalinput}[1]{\IfFileExists{\fiscaldir/#1}{\input{\fiscaldir/#1}}%
  {\figbox{run \texttt{reports/fiscal\_figures.py} to fill this table}{1.2cm}}}
\newcommand{\figbox}[2]{\fbox{\parbox[c][#2][c]{\dimexpr\linewidth-2\fboxsep-2\fboxrule}%
  {\centering\small\itshape #1}}}

\title{\vspace{-1.6cm}Balanced-growth calibration: Greece}
\author{}
\date{\small $g=\gval$,\; $n_0=\nval$,\; $n_\infty=\nTval$,\;
      $r_B=\rBval$,\; $\Gamma_0-1=\Gammaval$,\; $\Gamma_T-1=\GammaTval$}

\begin{document}
\maketitle
\vspace{-0.6cm}

\noindent\small
The model is a small open economy with overlapping generations, solved in
detrended per-capita units. A hat denotes $\hat x_t=X_t/(Z_tN_t)$, with
$Z_t=(1+g)^t$ labour productivity and $N_t$ the population, so stocks
accumulate at the gross rate $\Gamma_t=(1+g)(1+n_t)$. On the balanced growth
path $\hat x$ is constant, levels grow at $\Gamma-1$ and per-capita
quantities at $g$. Preferences are $\log c-\nu\ell^{1+\varphi}/(1+\varphi)$.
Households hold domestic capital and foreign assets at the world return $r$;
government debt is held by official creditors at the real rate $r_B$. TFP $A$
is set so that $\hat y=1$ in 2023, so ratios to output read as levels. The
period is a year. Wealth sits on a grid of 100 nodes from 0 to 50 times
output per person, denser near zero (spacing $u^{1.5}$). Saving choices are
grid nodes, and aggregates come from the distribution over household states
carried forward by age rather than from simulation.



\vspace{0.1cm}
\begin{table}[!htbp]\centering\footnotesize
\caption{Parameters}\vspace{-2pt}
\tabinput{params_body.tex}
\end{table}

\noindent\footnotesize
Sources for the parameters not taken straight from a published table.
$K_g/Y$ is the IMF ICSD public capital stock for 2019, carried to 2023 by
perpetual inventory with national-accounts public investment and real GDP.
$\delta_g$ is estimated from the same stock over 2000--2019 and is within
0.001 of $I_g/K_g-(\Gamma_0-1)$. $r_B$ is a real rate, the 2026--60 mean of
the real effective interest rate in the Commission's debt projection for
Greece, which deflates the nominal effective rate on the stock by GDP deflator
growth (Spring 2026 forecast, Debt Sustainability Monitor layout). The past
gives the same number: the implicit nominal rate averaged 1.9\% over 2012--24
and deflator growth 0.9\%. $b_{min}$ is the national pension, \EUR{413.76} a
month from January 2023, over output per person aged 25--84 (\EUR{224.7}bn of
GDP over 7.56 million people, \EUR{29{,}722} each). The amount is from the
legislation and the 2024 Ageing Report country fiche; Eurostat reports only
aggregate pension spending.

The retirement age is the average effective age at which people start an
old-age, early or disability pension, 63.8 in 2022 against a statutory 67. To
2070 it follows the Ageing Report's projection (country fiche, Table 4): 65.5
in 2030, 66.4 in 2040, 66.6 in 2050, 67.4 in 2060 and 67.9 in 2070, linear in
between. After 2070 it rises one for one with life expectancy at 65 (men and
women averaged, EUROPOP2023) at the three-yearly reviews of Law 4336/2015, and
reaches 70.4 in 2100. A cohort retires at the age in force when it gets there,
67.7 for the cohort entering in 2023 and 69.9 for the one entering in 2050.
The model is annual, so each cohort is split between the two whole ages around
its average, in the proportions that reproduce it.

The replacement rate in year $t$ is $\rho^{pens}$ times an index equal to one
in 2023. The index is the Ageing Report's average pension relative to GDP per
employed person, pension spending over GDP divided by pensioners over
employment (country fiche, Tables 6 and 10): 0.85 in 2030, 0.77 in 2040, 0.67
in 2050 and 0.62 from 2060, linear in between. It applies to every pension in
payment, $b_{min}$ included, and households know the path from entry, as they
know their cohort's retirement age and survival rates. Unemployment rates by
education are the 2023 rates for ages 25--64 (Eurostat
\texttt{lfsa\_urgaed}). Education shares are for the population aged 25--84 in
2023, from the labour force survey for ages 25--74 (\texttt{lfsa\_pgaed}) and
the 2021 census for ages 75--84 (\texttt{cens\_21ae\_r2}).

\begin{table}[!htbp]\centering\footnotesize
\caption{Moments}\vspace{-2pt}
\tabinput{moments_body.tex}
\caption*{\footnotesize The SMM objective is $\sum_i (m_i-d_i)^2/d_i^2$, so
every target counts by its relative deviation. Interest$/Y$ and $I_g/Y$ are
set in the calibration and not listed. The untargeted rows compare income
inequality, with one caveat. Eurostat measures equivalised \emph{household}
disposable income; the model's figure is individual income net of income tax
and contributions. Taxes take the model Gini from 0.357 to 0.336 and the
p90/p10 ratio from 4.77 to 4.32, and pooling within households would lower
both further, so the model figures are upper bounds.}
\end{table}

\begin{table}[!htbp]\centering\footnotesize
\caption{Calibrated parameters against their direct data counterparts}\vspace{-2pt}
\tabinput{implied_body.tex}
\end{table}

\noindent\footnotesize
$\tau^p$ and $\rho^{pens}$ have direct data counterparts too. They differ
because the model levies $\tau^p$ on the whole wage bill, with no contribution
ceiling, and applies $\rho^{pens}$ to career-average earnings capacity, which
excludes hours worked, rather than to the final wage.

\section*{Demographics}
\vspace{-0.1cm}

\noindent\footnotesize
The calibration starts from the measured 2023 age distribution rather than
a stationary population.\footnote{A stationary population with the same
dependency ratio would miss the large cohorts about to retire, which drive the
projected rise in the ratio.} Each cohort faces the survival rates of its own
birth year. From 2023, survival and the size of entering cohorts follow the
EUROPOP2023 baseline to 2100, in the baseline and in every experiment. After
2100 mortality stays at its 2100 level and the growth of entering cohorts goes
to zero by 2120; the age distribution is constant from 2179 on.

Because the 2023 population is not stationary, the economy is in transition
even with no policy change. Prices do not move, since the return is set abroad
and the wage follows from it. Ageing reaches households through the retirement
age, higher for later cohorts, and through the pension index, which lowers the
pension a given career earns. Figure~\ref{fig:demography} shows the inputs.
Output per person falls to 2051 (Figure~\ref{fig:aggregates}) and pensions
rise to 2050 (Figure~\ref{fig:fiscal}) as the share aged 40--64 shrinks and
the number of retired per non-retired person grows; both turn around once the
smaller cohorts reach old age and the retirement age keeps rising.

\begin{table}[htbp]\centering\footnotesize
\caption{Age structure}\vspace{-3pt}
\tabinput{age_body.tex}
\end{table}

\begin{figure}[htbp]\centering
\IfFileExists{\rundir/demography.pdf}
  {\includegraphics[width=\linewidth]{\rundir/demography.pdf}}
  {\figbox{Figure: demographics --- run \texttt{reports/baseline\_figures.py}.}{6cm}}
\caption{\label{fig:demography}Demographic inputs, 2023--2120. \emph{Top
left}: the population aged 25--84 and the cohort entering at 25, relative to
2023. The model adds people only at 25, so the entering cohort is set each
year to make the model's population aged 25--84 equal EUROPOP2023's. All net
migration therefore passes through it. It is negative until the 2050s, hence
the fall in 2024 and the movements through the 2030s. The drop in 2048 is
EUROPOP2023's own: its 2022 birth cohort is 9\% smaller than 2021's, matching
the recorded fall in live births from 85.3 to 75.9 thousand (Eurostat
\texttt{demo\_fasec}). The projection starts from the population of 1 January
2022 (the model's base year is 2023), so later cohorts are projected, and they
run above recorded births (79.2 and 77.8 thousand for 2023 and 2024 against
71.2 and 68.3), so the cohorts entering from 2049 on are 10--14\% larger than
recorded births imply. The dip in 2084--90 is the exit at 85 of the small
2024--29 cohorts. \emph{Top right}: age shares; the share aged 65--84 rises
from 0.26 to 0.41 by 2050 as the cohorts born in the 1960s and 1970s retire.
\emph{Bottom left}: persons aged 65--84 per person aged 25--64 peak at 0.69 in
2051, from 0.36 in 2023; the retired per non-retired person peak at 0.60 and
fall to 0.29 by 2100, below today's value, because the retirement age rises
and people above 84 are outside the model. \emph{Bottom right}: the average
effective retirement age, 64.0 in 2023 to 70.4 in 2100, and 65 plus life
expectancy at 65, 85.4 to 92.8 years.}
\end{figure}

\clearpage
\section*{Baseline transition}\vspace{-0.2cm}

\begin{figure}[htbp]\centering
\IfFileExists{\rundir/baseline_aggregates.pdf}
  {\includegraphics[width=\linewidth]{\rundir/baseline_aggregates.pdf}}
  {\figbox{Figure: aggregates --- run \texttt{reports/baseline\_figures.py}.}{4.5cm}}
\caption{\label{fig:aggregates}Baseline, no policy change. \emph{Left}:
detrended per-capita output, consumption and household wealth relative to
2023. Domestic capital and hours move with output, since the world return
fixes $K^{dom}/L$ and public capital per person is constant. \emph{Centre and
right}: ratios to output. Private investment,
$I_t=\Gamma_t K^{dom}_{t+1}-(1-\delta)K^{dom}_t$, is shown as a 5-year centred
moving average, because as the change in a stock some 15 times larger it
magnifies the stock's year-to-year movement about 20-fold. Household wealth is
$A=K^{dom}+NFA$. Everything here comes from the age structure and the rising
retirement age.}
\end{figure}

\begin{figure}[htbp]\centering
\IfFileExists{\rundir/baseline_fiscal.pdf}
  {\includegraphics[width=\linewidth]{\rundir/baseline_fiscal.pdf}}
  {\figbox{Figure: government accounts --- run \texttt{reports/baseline\_figures.py}.}{6cm}}
\caption{\label{fig:fiscal}Baseline, no policy change: government accounts
as shares of output. Tax rates, $G/Y$ and defence$/Y$ are fixed. Pensions
follow the index and the number of retirees; health spending and the tax
bases follow the age structure. Other net spending $O$ is the primary spending
less revenue that has no explicit line in the model, and no household pays or
receives it. It is set so that the primary balance equals the 2023 outturn
(1.95\% of output), the Commission's projection over 2024--60, and from 2061
the balance that holds the debt ratio at its 2060 value; the lower left panel
shows that balance and the balance before $O$. Debt is the end-of-year stock
over the year's output, $d_t=d_{t-1}(1+r_B)/(1+g_t)-pb_t+sfa_t$ from 1.64 in
2023, with $g_t=\Gamma_{t-1}\hat y_t/\hat y_{t-1}-1$. The stock-flow
adjustment $sfa$ reproduces the ratios of 2024 (154.2\%) and 2025 (146.1\%),
is the projection's own over 2026--60 (0.8--1.1\% of output a year to 2032,
the deferred interest on the official loans) and zero after. Debt falls to
136.9\% in 2030 and 129.1\% in 2040, rises to 134.6\% in 2050 and settles at
125.9\% from 2060; the projection reaches 103.4\% in 2060. The gap comes from
the constant real rate and from lower output growth
(Figure~\ref{fig:projection}).}
\end{figure}

\begin{figure}[htbp]\centering
\IfFileExists{\rundir/baseline_projection.pdf}
  {\includegraphics[width=\linewidth]{\rundir/baseline_projection.pdf}}
  {\figbox{Figure: projection --- run \texttt{reports/baseline\_figures.py}.}{6cm}}
\caption{\label{fig:projection}The baseline against the Commission's debt
projection for Greece (Spring 2026 forecast, 2025--60). \emph{Top left}:
end-of-year debt over output. The two part ways through the real rate, which
in the projection is negative until 2032, and through output growth.
\emph{Top right}: the primary balance, the projection's by construction to
2060, and other net spending. \emph{Bottom left}: public pensions over output,
against the 2024 Ageing Report (14.5\% of GDP in 2022, 14.0\% in 2050, 12.0\%
in 2070). The model's level is the national-accounts 16\% of 2023; the Report
uses the Ageing Working Group definition. \emph{Bottom right}: real output
growth, the baseline as a 5-year centred mean. The projection's growth is the
forecast to 2027 and potential growth after; the model has no cycle.}
\end{figure}

\begin{table}[htbp]\centering\footnotesize
\caption{Baseline and the Commission's projection, \% of output}\vspace{-2pt}
\label{tab:projection}
\tabinput{projection_body.tex}
\caption*{\footnotesize End-of-year debt over the year's output. The primary
balance is the projection's to 2060. $O/Y$ is other net spending, negative
where the revenue the model lacks exceeds the spending it lacks. Revenue,
pensions and public health are the model's own lines.}
\end{table}

\noindent\footnotesize
Pensions fall from 16.0\% of output in 2023 to 14.2\% in 2030, because the
index falls faster than the number of retirees grows, climb back to 15.9\% in
2050 as the large cohorts retire, and then fall to 13.2\% in 2060 and 10.3\%
in 2070 once the index settles and the retirement age keeps rising. The Ageing
Report has 14.5\% of GDP in 2022, 12.7\% in 2030, 14.0\% in 2050 and 12.0\% in
2070. Public health spending rises from 5.4\% to 7.5\% of output by 2050, with
the unit cost by age growing at the trend. Revenue rises from 35.0\% to
38.5\% of output by 2050, mostly consumption tax: households save more for
their lower pensions, wealth rises from 4.0 to 5.3 times output, and the
consumption it finances is taxed. The primary balance before $O$ is $-6.9$\%
of output in 2023, $-4.9$\% in 2030 and $-5.6$\% in 2050, so $O/Y$ goes from
$-8.9$\% in 2023 ($-11.2$\% in 2025) to $-7.2$\% in 2030, $-6.2$\% in 2040 and
$-5.4$\% in 2060, and then up to $+0.4$\% by 2100 as pensions keep falling; a
primary deficit of about 0.8\% of output holds debt at 126\%. Output per
person bottoms out at 0.87 of its 2023 level in 2051 and recovers to 1.06 by
2100 on the rising retirement age. Real growth averages 0.65\% a year over
2026--60, against 1.08\% in the projection. The gap is widest in 2041--50
(0.20\% against 1.17\%), when the retirement age rises only 0.2 years over the
decade while the working-age population shrinks.

\begin{table}[htbp]\centering\footnotesize
\caption{Growth of the detrended series on the balanced growth path}\vspace{-2pt}
\label{tab:growth}
\tabinput{growth_body.tex}
\end{table}

\noindent\footnotesize
Table~\ref{tab:growth} reports the mean log growth per year,
$\frac{1}{T-t_0}\sum_{t\ge t_0}\log(\hat x_{t+1}/\hat x_t)$, from
$t_0=2179$, when the age distribution becomes constant, to the end of the
horizon. With correct detrending every series is flat. A series left in levels
would show $\Gamma_T-1$ (1.7\%) and one left per capita $g$ (also 1.7\%). All
entries are zero to three decimals.

\clearpage
\section*{Policy experiments}\vspace{-0.2cm}

\noindent\footnotesize
Two permanent spending increases start in 2023, each worth 2\% of output.
Government consumption $G$ rises by 2\% of each year's output. Public
investment $I_g$ rises by 2\% of 2023 output and stays there in detrended
per-capita terms, because public capital enters production and an increase
set as a share of output would depend on its own effect. Neither enters
utility. Under debt financing the primary deficit is borrowed abroad at $r_B$
and no tax changes. Under the two labour-tax rules $\tau_l$ rises by a
constant $\Delta\tau_l$ from 2023, set so that in 2202, the last year of the
horizon, either the debt ratio or net foreign assets over output are back at
their baseline values. Other net spending keeps its baseline share of each
run's output and the stock-flow adjustment its baseline levels. The policy is
announced in 2023, and the wealth each cohort holds that year is the
baseline's. Net foreign assets are household foreign assets less government
debt. The cumulative multiplier is the sum of the output response over the
sum of the spending increase, both as aggregate flows over the horizon.

Debt-financed government consumption leaves output, consumption and hours
exactly where they were. It enters neither utility nor production, and the
debt is held abroad, so no household budget changes. The multiplier is zero,
and the debt ratio climbs to 207\% of output in 2060 and 339\% in 2200.
Bringing the ratio back to 126\% by 2202 takes $\Delta\tau_l=4.19$ percentage
points. Hours and output fall 0.8\% in 2023 and 0.4\% on average over
2023--32, and are back at the baseline by 2050; consumption is 3.2\% below the
baseline in 2023 and 4.8\% below in 2100, wealth 4.6\% below in 2100.
Returning net foreign assets to their baseline ratio instead takes 4.67
points and leaves the debt ratio at 102\% in 2200.

Public investment raises public capital per person from 0.70 of 2023 output
to 0.99 in 2050 and 1.05 in 2100. Debt-financed, it lowers output by 1.3\% in
2023, because households expect higher wages and work less now, and raises it
by 2.7\% in 2050 and 3.0\% in 2100. Consumption is 1.2\% above the baseline in
2023 and 3.1\% above in 2100. Output is below the baseline for three years and
above it from 2026; the cumulative multiplier is 0.20 over 2023--32 and 1.49
over the horizon. The debt ratio reaches 203\% in 2060 and 299\% in 2200. The
labour-tax rules need 3.42 points (debt target) and 3.82 points
(net-foreign-asset target). Output is 1.6\% below the baseline in 2023 under
both and 2.8--2.9\% above in 2050; consumption ends 0.9\% and 1.4\% below in
2100.

\begin{table}[htbp]\centering\footnotesize
\caption{Policy experiments: responses by financing rule}\vspace{-2pt}
\fiscalinput{fiscal_body.tex}
\caption*{\footnotesize End-of-year debt over the year's output; the
baseline has 126\% from 2060 on and net foreign assets of $-0.38$ of output
in 2200. The 2023--32 column is the mean deviation over those years.}
\end{table}

\begin{figure}[htbp]\centering
\IfFileExists{\fiscaldir/fiscal_G.pdf}
  {\includegraphics[width=\linewidth]{\fiscaldir/fiscal_G.pdf}}
  {\figbox{Figure: G shock --- run \texttt{reports/fiscal\_figures.py}.}{6cm}}
\caption{\label{fig:G}Government consumption up by 2\% of output from 2023, under the
three financing rules. \emph{Top}: output and consumption, \% deviation from
the baseline. \emph{Bottom}: debt at the end of the year and net foreign
assets, both over output, with the baseline.}
\end{figure}

\begin{figure}[htbp]\centering
\IfFileExists{\fiscaldir/fiscal_Ig.pdf}
  {\includegraphics[width=\linewidth]{\fiscaldir/fiscal_Ig.pdf}}
  {\figbox{Figure: I\_g shock --- run \texttt{reports/fiscal\_figures.py}.}{6cm}}
\caption{\label{fig:Ig}Public investment up by 2\% of 2023 output from 2023, under the
three financing rules; panels as in Figure~\ref{fig:G}.}
\end{figure}

\end{document}
